\chapter{Example Character Concepts}
\label{ch:example-concepts}
\index{character concepts!examples}

This chapter presents example character concepts to illustrate how the game's systems can create diverse and interesting heroes. These are \textbf{examples only}—not prescriptive templates or exhaustive lists. Use them for inspiration, as pre-generated characters, or as starting points for your own unique creations.

\section{Important Disclaimer}
\label{sec:concepts-disclaimer}
\index{character concepts!disclaimer}

\textbf{These examples are provided for illustrative purposes only.} They demonstrate how the game's mechanics can support different character archetypes and play styles. You are encouraged to:
\begin{itemize}
\item Modify these concepts to fit your preferences
\item Create completely original characters
\item Mix and match elements from different examples
\item Work with your Game Master to develop unique concepts
\end{itemize}

The game system is designed to support a wide variety of character types beyond these examples.

\section{How to Use These Examples}
\label{sec:concepts-usage}
\index{character concepts!usage}

Each concept includes:
\begin{itemize}
\item \textbf{Concept Overview}: Narrative identity and role
\item \textbf{Traditional Analogues}: Common RPG classes that share similar themes — \textit{not} mechanical restrictions, just touchstones
\item \textbf{Mechanical Foundation}: Suggested starting capabilities
\item \textbf{Play Style}: How the character typically engages with challenges
\item \textbf{Development Path (by Tier)}: What to buy at Tier I, II, and III (see \ref{sec:tiers})
\item \textbf{Tactical Advice}: Tips for playing the role effectively
\item \textbf{Story Hooks}: Plot opportunities for the Game Master
\item \textbf{Legacy Note}: How this character might pass the torch (see \ref{ch:legacy-engine})
\item \textbf{Build Blocks}: A \emph{30 XP} starting build, plus an optional \emph{34 XP} variant using Bonds/Complications (+4 XP) (see \ref{ch:advancement} for XP rules)
\end{itemize}

\section{1. The Guardian (Fighter / Paladin / Knight)}
\label{sec:concept-guardian}
\index{character concepts!guardian}

\textbf{Concept}: A protector who stands between danger and those they've sworn to defend. \emph{Steel in hand, vow in heart.} Whether a chivalrous knight, a grizzled mercenary, or a holy warrior of the Everflame, the Guardian makes others safe through sacrifice.

\textbf{Traditional Analogues}: Fighter, Paladin, Barbarian (defensive), Knight, Samurai, Legionary

\textbf{Mechanical Foundation}:
\begin{itemize}
\item \textbf{Primary}: Body (strength, endurance), Spirit (willpower, resolve)
\item \textbf{Skills}: Melee, Athletics, Command, Endurance
\item \textbf{Talents}: Defensive Survival, Combat Reflexes, Conditioning (see \ref{ch:talents})
\end{itemize}

\textbf{Play Style}:
\begin{itemize}
\item Frontline combat and protection
\item Drawing attention away from allies (taunts, positioning)
\item Using presence and authority to control situations
\item Taking risks to protect others — absorbing Harm meant for teammates
\end{itemize}

\paragraph{Development Path (by Tier)}:
\begin{itemize}
\item \textbf{Tier I (0–40 XP)}: Raise Body to 3, Melee to 2, take \emph{Defensive Survival} and \emph{Combat Reflexes}. Use a shield.
\item \textbf{Tier II (41–90 XP)}: Raise Body to 4, take \emph{Conditioning} and \emph{Tactical Movement}. Consider a War Mount Asset (8 XP) for mobility.
\item \textbf{Tier III (91–150 XP)}: Take \emph{Battlefield Mastery} or \emph{Mystic Bulwark}. Upgrade armor to Heavy if you have Strength. Start training a successor (Cap 3+ follower).
\end{itemize}

\paragraph{Tactical Advice}:
\begin{itemize}
\item \textbf{Use the Defend action} when outnumbered; it improves your Position by one step and can protect allies.
\item \textbf{Spend Boons on re‑rolling failed defense} rolls, not on offense. Your job is to survive.
\item \textbf{Position matters more than damage.} A Dominant position lets you re‑roll failures – use shoves, terrain, and cover to stay in control.
\item \textbf{When you take Harm, remember it clears Fatigue.} Sometimes a small wound is worth the reset.
\end{itemize}

\textbf{Story Hooks}:
\begin{itemize}
\item Who or what are they protecting? A person, a place, an ideal?
\item What oath or duty drives them — sworn, inherited, or self‑imposed?
\item What happens if they fail in their protection?
\item What personal costs do they bear for their role (scars, lost loved ones, broken oaths)?
\end{itemize}

\textbf{Legacy Note}: A loyal follower (squire, apprentice) can become your successor, inheriting your sworn duty and reputation. See \ref{ch:legacy-engine}.

\paragraph{Build Blocks.}
\textbf{Starting Build (30 XP).}
\begin{itemize}
\item \textbf{Attributes} (Cost = rating $\times$ 3 XP): Body 3 (9), Spirit 2 (6), Wits 1 (3), Presence 1 (3) $\rightarrow$ \textbf{21 XP}
\item \textbf{Skills} (Cost = level $\times$ 2 XP): Melee 2 (4), Athletics 1 (2), Command 1 (2) $\rightarrow$ \textbf{8 XP}
\item \textbf{Total}: 29 XP (bank 1 XP)
\end{itemize}
\textbf{With Bonds/Complications (34 XP).}
\begin{itemize}
\item Add \textbf{Talent}: Combat Reflexes (2 XP)
\item Add \textbf{Talent}: Defensive Survival (3 XP)
\item \textbf{Total}: 29 + 5 = 34 XP
\end{itemize}

\section{2. The Mage (Wizard / Sorcerer / Arcanist)}
\label{sec:concept-mage}
\index{character concepts!mage}

\textbf{Concept}: A wielder of arcane power who reshapes reality through study, intuition, or pact. \emph{Candlesmoke, marginalia, and dangerous truths.} From the rational Wizard of Thepyrgos to the instinctive Sorcerer of the Bloodlands, magic flows through them — and sometimes through them.

\textbf{Traditional Analogues}: Wizard, Sorcerer, Warlock (pact magic), Arcanist, Runekeeper

\textbf{Mechanical Foundation}:
\begin{itemize}
\item \textbf{Primary}: Wits (theory, precision), Spirit (willpower, channeling)
\item \textbf{Skills}: Lore, Arcana, Investigation, (optionally) Sway (for social casting)
\item \textbf{Talents}: Spellcraft (freeform), Codex (runekeeper), Patron’s Symbol (invoker)
\end{itemize}

\textbf{Play Style}:
\begin{itemize}
\item Information gathering and magical analysis
\item Solving puzzles and mysteries with arcane insight
\item Using magic to gain advantages in combat and social scenes
\item Researching new spells and rituals between adventures
\item Managing Obligation (for Runekeepers) or Backlash (for free casters)
\end{itemize}

\paragraph{Development Path (by Tier)}:
\begin{itemize}
\item \textbf{Tier I}: Take \emph{Spellcraft} (6 XP) or \emph{Codex} + \emph{Thiasos} (6 XP). Raise Wits to 3, Arcana/Lore to 2.
\item \textbf{Tier II}: Add \emph{Elemental Focus} (4 XP) or \emph{Ritual Mastery} (4 XP). Raise Spirit to 3 for more Obligation capacity.
\item \textbf{Tier III}: Take \emph{Arcane Dominance} or \emph{Crack Specialist} (if Invoker). Build a library (Standard Asset, 8 XP) to research faster.
\end{itemize}

\paragraph{Tactical Advice}:
\begin{itemize}
\item \textbf{Push It} only when necessary – extra Obligation or Fatigue can snowball. Use Boons to improve Position before rolling.
\item \textbf{Free casters:} Use cantrips (single tag) for safe, repeatable effects. Reserve multi‑tag spells for critical moments.
\item \textbf{Runekeepers:} Your Patron’s Gift gives +1 Melee and +1 Thematic. Use it even if you aren’t a fighter – the Thematic skill can be Sway, Lore, or Stealth.
\item \textbf{Invokers:} \emph{Crack the Seal} is for emergencies. Restore your Symbol during downtime before you need it again.
\end{itemize}

\textbf{Story Hooks}:
\begin{itemize}
\item What forbidden knowledge are they seeking?
\item How did they acquire their magical talent — bloodline, study, pact, accident?
\item What dangerous secrets might they uncover?
\item Who opposes their research (rival mages, witch‑hunters, the Church)?
\end{itemize}

\textbf{Legacy Note}: A trusted apprentice (follower) can inherit your Codex or Symbol collection, continuing your magical legacy. See \ref{ch:legacy-engine}.

\paragraph{Build Blocks.}
\textbf{Starting Build (30 XP).}
\begin{itemize}
\item \textbf{Attributes}: Wits 3 (9), Spirit 2 (6), Body 1 (3), Presence 1 (3) $\rightarrow$ \textbf{21 XP}
\item \textbf{Skills}: Lore 2 (4), Arcana 1 (2) $\rightarrow$ \textbf{6 XP}
\item \textbf{Total}: 27 XP (bank 3 XP)
\end{itemize}
\textbf{With Bonds/Complications (34 XP).}
\begin{itemize}
\item Add \textbf{Talent}: Spellcraft (6 XP) using banked 3 + 4 = 7 XP; bank 1 XP
\item \textbf{Total}: 27 + 6 = 33 XP (bank 1 XP)
\end{itemize}

\section{3. The Rogue (Thief / Assassin / Scout)}
\label{sec:concept-rogue}
\index{character concepts!rogue}

\textbf{Concept}: A shadow who moves where others cannot, taking what others cannot keep. \emph{Quiet footfalls, hawk eyes, and no loyalties but coin.} From the street‑wise cutpurse of Silkstrand to the deadly assassin of Zakov, the Rogue survives by wits, speed, and one hidden knife.

\textbf{Traditional Analogues}: Rogue, Thief, Assassin, Scout, Ranger (stealth focus), Spy

\textbf{Mechanical Foundation}:
\begin{itemize}
\item \textbf{Primary}: Wits (perception, quick thinking), Body (agility, stealth)
\item \textbf{Skills}: Stealth, Subterfuge, Investigation (or Streetwise), (optionally) Melee (light weapons)
\item \textbf{Talents}: Backstab, Elusive Dodge, Shadow Dance (later – see \ref{ch:talents})
\end{itemize}

\textbf{Play Style}:
\begin{itemize}
\item Scouting ahead and gathering intelligence
\item Infiltrating secured areas and disabling traps
\item Ambush and skirmish tactics — striking from stealth
\item Finding paths, contacts, and secrets in cities and wilderness
\end{itemize}

\paragraph{Development Path (by Tier)}:
\begin{itemize}
\item \textbf{Tier I}: Raise Wits to 3, Stealth/Subterfuge to 2. Take \emph{Backstab} (6 XP) – you can afford it at 34 XP start.
\item \textbf{Tier II}: Take \emph{Elusive Dodge} (6 XP). Raise Body to 3 for better tracking. Consider a Minor Asset (Hidden Cache, 4 XP).
\item \textbf{Tier III}: Take \emph{Shadow Dance} (10 XP) and \emph{Deathblow} (12 XP) – the classic assassination chain. Build a spy network (Standard Asset, 8 XP).
\end{itemize}

\paragraph{Tactical Advice}:
\begin{itemize}
\item \textbf{Always aim for Dominant Position before attacking from stealth.} Spend a Boon to improve Position if needed.
\item \textbf{Use the environment.} Ask the GM: “Is there a chandelier? A shadowed alcove? Loose gravel?” Use it to gain advantage.
\item \textbf{Backstab + Shadow Dance + Deathblow} is lethal, but requires Boons. Engage only when you have 2–3 Boons saved.
\item \textbf{When outmatched, retreat.} Spend a Boon to make a Controlled withdrawal instead of Desperate.
\end{itemize}

\textbf{Story Hooks}:
\begin{itemize}
\item What crime or secret forced them into the shadows?
\item What treasure or truth are they hunting?
\item How do they balance loyalty to allies with self‑preservation?
\item Who wants them dead — and why?
\end{itemize}

\textbf{Legacy Note}: A young pickpocket or informant (follower) could inherit your network of fences and contacts. See \ref{ch:legacy-engine}.

\paragraph{Build Blocks.}
\textbf{Starting Build (30 XP).}
\begin{itemize}
\item \textbf{Attributes}: Wits 3 (9), Body 2 (6), Presence 1 (3), Spirit 1 (3) $\rightarrow$ \textbf{21 XP}
\item \textbf{Skills}: Stealth 2 (4), Subterfuge 2 (4) $\rightarrow$ \textbf{8 XP}
\item \textbf{Total}: 29 XP (bank 1 XP)
\end{itemize}
\textbf{With Bonds/Complications (34 XP).}
\begin{itemize}
\item Add \textbf{Talent}: Backstab (6 XP) – needs 6 XP; bank 1 + 4 = 5 XP → not enough. Alternative: Add \textbf{Skill}: Perception 1 (2 XP) and \textbf{Talent}: Combat Reflexes (2 XP) and bank 2 XP. Or take a Minor Asset: Hidden Cache (4 XP) and bank 2 XP.
\item \textbf{Suggested}: Hidden Cache + bank 2 XP → total 33 XP.
\end{itemize}

\section{4. The Cleric (Priest / Healer / Exorcist)}
\label{sec:concept-cleric}
\index{character concepts!cleric}

\textbf{Concept}: A vessel of divine or spiritual power, channeling healing, truth, and judgment. \emph{Sacred flame, holy water, and a prayer on the lips.} From the Lampers of Ecktoria to the fog‑named exorcists of the Mistlands, the Cleric binds wounds and banishes darkness.

\textbf{Traditional Analogues}: Cleric, Priest, Paladin (healing focus), Exorcist, Shaman, Healer

\textbf{Mechanical Foundation}:
\begin{itemize}
\item \textbf{Primary}: Spirit (faith, willpower), Presence (authority, channeling)
\item \textbf{Skills}: Lore (theology, rites), Medicine, Sway (prayer, exhortation), (optionally) Melee (symbolic weapon)
\item \textbf{Talents}: Invoker (Patron’s Symbol) or Runekeeper (Thiasos + Codex of a divine Patron like Everflame, Oath of Flame \& Light, Pale Shepherd – see \ref{ch:magic})
\end{itemize}

\textbf{Play Style}:
\begin{itemize}
\item Healing and protecting allies
\item Purifying curses, diseases, and undead
\item Ley line and ritual magic (wards, consecrations)
\item Using authority of faith to command or persuade
\end{itemize}

\paragraph{Development Path (by Tier)}:
\begin{itemize}
\item \textbf{Tier I}: Raise Spirit to 3, Lore/Medicine to 2. Take a divine Patron’s \emph{Symbol} (4 XP) or \emph{Codex} (4 XP).
\item \textbf{Tier II}: Take \emph{Symbol Resonance} (4 XP) or \emph{Efficient Invocation} (4 XP). Raise Presence to 3 for better social influence.
\item \textbf{Tier III}: Take \emph{Crack Specialist} or \emph{Patron’s Favor}. Build a hospice or shrine (Standard Asset, 8 XP) to support your community.
\end{itemize}

\paragraph{Tactical Advice}:
\begin{itemize}
\item \textbf{Healing is powerful but slow.} In combat, stabilise Harm 2+ allies first; use Mend (Fatigue recovery) on everyone else.
\item \textbf{Your Patron’s Symbol grants a free +1 to ritual rolls.} Keep it displayed.
\item \textbf{Consecrate ground before a big fight} – \emph{Rite of Consecrated Ground} can turn the tide against undead or demons.
\item \textbf{Obligation is story fuel.} When your Patron calls, answer – it will lead to adventure.
\end{itemize}

\textbf{Story Hooks}:
\begin{itemize}
\item What crisis of faith have they survived?
\item What heresy or supernatural threat haunts them?
\item How do they serve both the living and the dead?
\item What price do they pay for divine power?
\end{itemize}

\textbf{Legacy Note}: An acolyte or lay healer (follower) can take up your Patron’s Symbol and continue your mission. See \ref{ch:legacy-engine}.

\paragraph{Build Blocks.}
\textbf{Starting Build (30 XP).}
\begin{itemize}
\item \textbf{Attributes}: Spirit 3 (9), Presence 2 (6), Wits 1 (3), Body 1 (3) $\rightarrow$ \textbf{21 XP}
\item \textbf{Skills}: Medicine 2 (4), Lore 1 (2), Sway 1 (2) $\rightarrow$ \textbf{8 XP}
\item \textbf{Total}: 29 XP (bank 1 XP)
\end{itemize}
\textbf{With Bonds/Complications (34 XP).}
\begin{itemize}
\item Add \textbf{Talent}: Patron’s Symbol (4 XP) for a divine Patron (e.g., Oath of Flame \& Light), using banked 1 + 4 = 5 XP; bank 1 XP
\item \textbf{Total}: 33 XP (bank 1 XP)
\end{itemize}

\section{5. The Bard (Diplomat / Skald / Courtier)}
\label{sec:concept-bard}
\index{character concepts!bard}

\textbf{Concept}: A master of words, music, and social influence who changes minds and opens doors. \emph{A smile for the foyer, steel for the parlor.} From the Skald of Linn to the Courtier of Vhasia, the Bard wields charm, rumor, and reputation as weapons.

\textbf{Traditional Analogues}: Bard, Diplomat, Courtier, Skald, Spy (social), Negotiator

\textbf{Mechanical Foundation}:
\begin{itemize}
\item \textbf{Primary}: Presence (charm, performance), Wits (reading people, timing)
\item \textbf{Skills}: Sway, Performance, Insight, (optionally) Lore (customs, history)
\item \textbf{Talents}: Silver Tongue, Inspire, Network Builder (or Read Emotions)
\end{itemize}

\textbf{Play Style}:
\begin{itemize}
\item Social interaction and negotiation
\item Gathering information through rumors and contacts
\item Resolving conflicts without violence (or making violence fashionable)
\item Building alliances, manipulating factions
\end{itemize}

\paragraph{Development Path (by Tier)}:
\begin{itemize}
\item \textbf{Tier I}: Raise Presence to 3, Sway to 2. Take \emph{Silver Tongue} (2 XP) and \emph{Inspire} (3 XP).
\item \textbf{Tier II}: Take \emph{Command Presence} (4 XP) and \emph{Network Builder} (4 XP). Raise Wits to 3 for social reading.
\item \textbf{Tier III}: Take \emph{High Stance Duelist} (if you use a blade) or \emph{Battlefield Terror}. Build a spy network (Standard Asset, 8 XP) or a courtly title (Major Asset, 12 XP).
\end{itemize}

\paragraph{Tactical Advice}:
\begin{itemize}
\item \textbf{Social scenes are your domain.} Use \emph{Inspire} to give allies +1 die or +1 Boon – it can turn a negotiation.
\item \textbf{Flashback declarations} (spend 1 Boon) are perfect for bards. “I already bribed the guard last night.”
\item \textbf{Don’t neglect combat skills entirely.} A rapier and \emph{High Stance Duelist} can make you a credible threat.
\item \textbf{Use \emph{Intricate} descriptions} to gain free re‑rolls on 1s. Describe your performance, your charming smile, your family crest.
\end{itemize}

\textbf{Story Hooks}:
\begin{itemize}
\item What major conflict are they trying to resolve?
\item What secret do they know that could topple a dynasty?
\item How do they handle betrayal or failed negotiations?
\item What personal relationships affect their diplomacy?
\end{itemize}

\textbf{Legacy Note}: A protégé or herald (follower) can inherit your reputation and contacts, keeping your network alive. See \ref{ch:legacy-engine}.

\paragraph{Build Blocks.}
\textbf{Starting Build (30 XP).}
\begin{itemize}
\item \textbf{Attributes}: Presence 3 (9), Wits 2 (6), Spirit 1 (3), Body 1 (3) $\rightarrow$ \textbf{21 XP}
\item \textbf{Skills}: Sway 2 (4), Performance 1 (2), Insight 1 (2) $\rightarrow$ \textbf{8 XP}
\item \textbf{Total}: 29 XP (bank 1 XP)
\end{itemize}
\textbf{With Bonds/Complications (34 XP).}
\begin{itemize}
\item Add \textbf{Talent}: Inspire (3 XP) — once/scene, gain +1 die and potentially +1 Boon for bonded ally, using banked 1 + 4 = 5 XP; bank 2 XP
\item \textbf{Total}: 32 XP (bank 2 XP)
\end{itemize}

\section{6. The Ranger (Hunter / Tracker / Beastmaster)}
\label{sec:concept-ranger}
\index{character concepts!ranger}

\textbf{Concept}: A wilderness expert who navigates dangerous territories and hunts down quarry. \emph{Quiet footfalls, hawk eyes, and the long road.} From the Ykrul scout to the Valewood warden, the Ranger lives where roads end.

\textbf{Traditional Analogues}: Ranger, Hunter, Tracker, Beastmaster, Pathfinder, Wilderness Scout

\textbf{Mechanical Foundation}:
\begin{itemize}
\item \textbf{Primary}: Wits (perception, navigation), Body (agility, endurance)
\item \textbf{Skills}: Survival, Stealth, Nature (or Animal Handling), Archery or Melee
\item \textbf{Talents}: Wilderness Lore, Keen Senses, (optionally) Animal Companion
\end{itemize}

\textbf{Play Style}:
\begin{itemize}
\item Scouting ahead and gathering intelligence
\item Wilderness survival and navigation
\item Tracking creatures or people across terrain
\item Ambushing from concealment, using terrain
\item (Optional) Bonding with an animal companion
\end{itemize}

\paragraph{Development Path (by Tier)}:
\begin{itemize}
\item \textbf{Tier I}: Raise Wits to 3, Survival/Stealth to 2. Take \emph{Wilderness Lore} (2 XP) and \emph{Keen Senses} (2 XP).
\item \textbf{Tier II}: Take \emph{Trackless Step} (2 XP) and \emph{Elusive Dodge} (6 XP) or a War Hound (Minor Asset, 4 XP). Raise Body to 3.
\item \textbf{Tier III}: Take \emph{Shadow Dance} (10 XP) (if melee) or \emph{Relentless} (6 XP) for traversal. Upgrade your animal companion to a War Mount (Standard Asset, 8 XP) if allowed.
\end{itemize}

\paragraph{Tactical Advice}:
\begin{itemize}
\item \textbf{Use terrain to your advantage.} High ground, cover, and difficult terrain are your best friends. Ask the GM to set Position based on these.
\item \textbf{Scout ahead} even when not required. Knowledge of what’s coming negates ambushes and reduces travel clocks.
\item \textbf{If you have an animal companion}, use it to harass enemies (Distract \& Draw independent action) rather than direct combat.
\item \textbf{Survival 3+ means you rarely need to roll for basic tasks.} Leverage that for guaranteed successes in wilderness travel.
\end{itemize}

\textbf{Story Hooks}:
\begin{itemize}
\item What uncharted territory are they exploring?
\item What secrets have they discovered in the wild?
\item How do they balance civilization and wilderness?
\item What prey (beast or man) still eludes them?
\end{itemize}

\textbf{Legacy Note}: Your animal companion or a young tracker (follower) could carry on your legacy. See \ref{ch:legacy-engine}.

\paragraph{Build Blocks.}
\textbf{Starting Build (30 XP).}
\begin{itemize}
\item \textbf{Attributes}: Wits 3 (9), Body 2 (6), Spirit 1 (3), Presence 1 (3) $\rightarrow$ \textbf{21 XP}
\item \textbf{Skills}: Survival 2 (4), Stealth 2 (4) $\rightarrow$ \textbf{8 XP}
\item \textbf{Total}: 29 XP (bank 1 XP)
\end{itemize}
\textbf{With Bonds/Complications (34 XP).}
\begin{itemize}
\item Add \textbf{Talent}: Keen Senses (2 XP)
\item Add \textbf{Talent}: Wilderness Lore (2 XP)
\item \textbf{Total}: 29 + 4 = 33 XP (bank 1 XP)
\end{itemize}

\section{7. The Survivor (Veteran / Outcast / Hardened)}
\label{sec:concept-survivor}
\index{character concepts!survivor}

\textbf{Concept}: Someone who has endured hardship and developed resilience. \emph{Scars are maps; read them well.} From the war‑torn mercenary to the plague survivor, the Survivor won’t break — they’ve already been broken and rebuilt.

\textbf{Traditional Analogues}: Barbarian (resilience focus), Veteran, Ranger (endurance), Drifter, Refugee

\textbf{Mechanical Foundation}:
\begin{itemize}
\item \textbf{Primary}: Spirit (willpower, resolve), Body (toughness)
\item \textbf{Skills}: Endurance, Survival, (optionally) Perception (danger sense), Medicine (field)
\item \textbf{Talents}: Iron Stomach, Conditioning, Second Wind (see \ref{ch:talents})
\end{itemize}

\textbf{Play Style}:
\begin{itemize}
\item Enduring difficult conditions (lack of food, exposure, poison)
\item Overcoming physical and mental challenges through grit
\item Using experience to avoid dangers
\item Helping others survive hardships
\item Taking hits that would fell others and still standing
\end{itemize}

\paragraph{Development Path (by Tier)}:
\begin{itemize}
\item \textbf{Tier I}: Raise Spirit to 3, Endurance/Survival to 2. Take \emph{Iron Stomach} (3 XP) and \emph{Second Wind} (2 XP).
\item \textbf{Tier II}: Take \emph{Conditioning} (4 XP) – your Body counts as +1 for Fatigue track. Raise Body to 3.
\item \textbf{Tier III}: Take \emph{Hollow-Touched} (8 XP) or \emph{Stone-Bound Interdictor} (8 XP) to convert Harm to Fatigue. Build a safehouse (Minor Asset, 4 XP) for respite.
\end{itemize}

\paragraph{Tactical Advice}:
\begin{itemize}
\item \textbf{Embrace the roller‑coaster effect.} Taking Harm clears Fatigue – sometimes a small wound is a tactical reset.
\item \textbf{Iron Stomach makes you immune to mundane poisons and spoiled food.} Volunteer to test unknown rations.
\item \textbf{Use \emph{Second Wind} to clear Fatigue before it overflows into Harm.}
\item \textbf{When outmatched, use the environment.} Collapse a wall, start a fire, or trigger a rockfall. Your survival skills tell you how.
\end{itemize}

\textbf{Story Hooks}:
\begin{itemize}
\item What trauma or hardship have they survived?
\item How has their past shaped their present?
\item What survival skills have saved them repeatedly?
\item How do they help others facing similar challenges?
\end{itemize}

\textbf{Legacy Note}: A fellow survivor (follower) can inherit your network of caches and safehouses. See \ref{ch:legacy-engine}.

\paragraph{Build Blocks.}
\textbf{Starting Build (30 XP).}
\begin{itemize}
\item \textbf{Attributes}: Spirit 3 (9), Body 2 (6), Wits 1 (3), Presence 1 (3) $\rightarrow$ \textbf{21 XP}
\item \textbf{Skills}: Endurance 2 (4), Survival 2 (4) $\rightarrow$ \textbf{8 XP}
\item \textbf{Total}: 29 XP (bank 1 XP)
\end{itemize}
\textbf{With Bonds/Complications (34 XP).}
\begin{itemize}
\item Add \textbf{Talent}: Iron Stomach (3 XP)
\item \textbf{Total}: 29 + 3 = 32 XP (bank 2 XP)
\end{itemize}

\section{8. The Innovator (Artificer / Engineer / Tinker)}
\label{sec:concept-innovator}
\index{character concepts!innovator}

\textbf{Concept}: A creative problem-solver who finds new solutions through craft and invention. \emph{Blueprints on napkins, tomorrow in your pocket.} From the Aeler engineer to the Gnome clockmaker, the Innovator builds what others only dream.

\textbf{Traditional Analogues}: Artificer, Engineer, Tinker, Alchemist, Inventor, Savant

\textbf{Mechanical Foundation}:
\begin{itemize}
\item \textbf{Primary}: Wits (design, analysis), Body (fine manipulation)
\item \textbf{Skills}: Craft, Tinker, Mechanics, (optionally) Alchemy, Investigation
\item \textbf{Talents}: Quick Study, Technical Expert, (optionally) Magical Laboratory asset (see \ref{ch:assets-followers})
\end{itemize}

\textbf{Play Style}:
\begin{itemize}
\item Solving problems with unique inventions
\item Creating or repairing specialized items
\item Providing services others cannot (lockpicking, trap disarm, device creation)
\item Using knowledge to gain advantages in planning
\end{itemize}

\paragraph{Development Path (by Tier)}:
\begin{itemize}
\item \textbf{Tier I}: Raise Wits to 3, Craft/Tinker to 2. Take \emph{Quick Study} (3 XP) and \emph{Workshop} (Minor Asset, 4 XP).
\item \textbf{Tier II}: Take \emph{Efficient Invocation} (if magical) or \emph{Weapon Mastery} (if martial). Raise Body to 3 for fine manipulation.
\item \textbf{Tier III}: Take \emph{Technical Expert} (custom) or \emph{Crack Specialist} (for magical items). Build a Guild Workshop (Standard Asset, 8 XP) to attract apprentices.
\end{itemize}

\paragraph{Tactical Advice}:
\begin{itemize}
\item \textbf{Preparation is your superpower.} Spend downtime building tools, traps, and contingencies. Your GM will reward that.
\item \textbf{Use \emph{Quick Study} to learn enemy weaknesses} on the fly. “I notice the golem’s joints are unarmoured – aim there.”
\item \textbf{Your workshop asset can jury‑rig solutions.} Spend 1 Boon to “have the right tool” for a lock, a trap, or a repair.
\item \textbf{When combat starts, use gadgets.} Caltrops, smoke bombs, and flash pellets can shift Position without a spell slot.
\end{itemize}

\textbf{Story Hooks}:
\begin{itemize}
\item What makes their specialty unique or valuable?
\item How did they acquire their special skills (guild, apprenticeship, self‑taught)?
\item What problems require their specific expertise?
\item Who seeks to control or exploit their abilities?
\end{itemize}

\textbf{Legacy Note}: An apprentice (follower) can inherit your workshop and blueprints, continuing your inventive tradition. See \ref{ch:legacy-engine}.

\paragraph{Build Blocks.}
\textbf{Starting Build (30 XP).}
\begin{itemize}
\item \textbf{Attributes}: Wits 3 (9), Body 2 (6), Presence 1 (3), Spirit 1 (3) $\rightarrow$ \textbf{21 XP}
\item \textbf{Skills}: Craft 2 (4), Tinker 2 (4) $\rightarrow$ \textbf{8 XP}
\item \textbf{Total}: 29 XP (bank 1 XP)
\end{itemize}
\textbf{With Bonds/Complications (34 XP).}
\begin{itemize}
\item Add \textbf{Talent}: Quick Study (3 XP) — +1 die to learn or adapt from observation, using banked 1 + 4 = 5 XP; bank 2 XP
\item \textbf{Total}: 32 XP (bank 2 XP)
\end{itemize}

\section{9. The Summoner (Pact‑Whisperer / Necromancer / Shaman)}
\label{sec:concept-summoner}
\index{character concepts!summoner}

\textbf{Concept}: A binder of spirits, elementals, and otherworldly entities. \emph{Circles of salt, whispered names, and bargains struck in shadows.} From the necromancer of the Mistlands to the shaman of the Steppe, the Summoner commands powers that others fear.

\textbf{Traditional Analogues}: Necromancer, Conjurer, Summoner, Shaman, Spirit Caller, Demonologist

\textbf{Mechanical Foundation}:
\begin{itemize}
\item \textbf{Primary}: Spirit (willpower, binding), Wits (rituals, names)
\item \textbf{Skills}: Arcana, Lore (spirits), (optionally) Command, Survival
\item \textbf{Talents}: Pact‑Whisperer (Lesser/Greater), Spirit Link, (optionally) Codex of Names (see \ref{ch:magic})
\end{itemize}

\textbf{Play Style}:
\begin{itemize}
\item Summoning creatures to fight, scout, or work
\item Binding entities to solve problems (guard, research, transport)
\item Negotiating with spirits and outsiders (social challenges)
\item Managing Leash clocks and Boon Finesse
\end{itemize}

\paragraph{Development Path (by Tier)}:
\begin{itemize}
\item \textbf{Tier I}: Take \emph{Pact‑Whisperer} (2 XP) and \emph{Lesser Pactwright} (2 XP). Raise Spirit to 3, Arcana/Lore to 2.
\item \textbf{Tier II}: Take \emph{Greater Pactwright} (4 XP) to summon Cap 3 spirits. Take \emph{Spirit Synergy} (4 XP) if you want two spirits.
\item \textbf{Tier III}: Take \emph{Spirit Link} (10 XP) or \emph{True Name Keeper} (15 XP). Build a Spirit Circle (Minor Asset, 4 XP) to reduce Leash ticks.
\end{itemize}

\paragraph{Tactical Advice}:
\begin{itemize}
\item \textbf{Boon Finesse} is your most important tool. Spend 1 Boon to clear 1 Leash tick – keep a reserve of 2–3 Boons for emergencies.
\item \textbf{Spirits are disposable but not free.} Use them for dangerous tasks (scouting a trapped room, triggering a trap) before risking your own skin.
\item \textbf{When the Leash fills, the spirit acts once to its nature.} Plan for it – position it near enemies before it breaks free.
\item \textbf{Negotiate with spirits.} They have desires – a promise of freedom, a drop of blood, a secret. A willing spirit ticks Leash slower.
\end{itemize}

\textbf{Story Hooks}:
\begin{itemize}
\item What entity are they bound to or seeking?
\item What price did they pay for their first summoning?
\item How do others react to their “dark” magic?
\item Who hunts them for their forbidden practices?
\end{itemize}

\textbf{Legacy Note}: A bonded spirit or an apprentice medium (follower) could inherit your spirit contracts, continuing your work. See \ref{ch:legacy-engine}.

\paragraph{Build Blocks.}
\textbf{Starting Build (30 XP).}
\begin{itemize}
\item \textbf{Attributes}: Spirit 3 (9), Wits 2 (6), Body 1 (3), Presence 1 (3) $\rightarrow$ \textbf{21 XP}
\item \textbf{Skills}: Arcana 2 (4), Lore 1 (2) $\rightarrow$ \textbf{6 XP}
\item \textbf{Total}: 27 XP (bank 3 XP)
\end{itemize}
\textbf{With Bonds/Complications (34 XP).}
\begin{itemize}
\item Add \textbf{Talent}: Pact‑Whisperer (2 XP)
\item Add \textbf{Talent}: Lesser Pactwright (2 XP)
\item \textbf{Total}: 27 + 4 = 31 XP (bank 3 XP; can add Familiar (2 XP) later or upgrade to Greater Pactwright)
\end{itemize}

\section{10. The Commander (Warlord / Tactician / Leader)}
\label{sec:concept-commander}
\index{character concepts!commander}

\textbf{Concept}: A leader who directs allies, controls battlefields, and shapes conflicts through strategy. \emph{Orders in a whisper, victory in formation.} From the Condotta captain to the Ekstorian legate, the Commander wins before blades are drawn.

\textbf{Traditional Analogues}: Warlord, Marshal, Tactician, Commander, Strategist, Banneret

\textbf{Mechanical Foundation}:
\begin{itemize}
\item \textbf{Primary}: Presence (authority, command), Wits (tactics, observation)
\item \textbf{Skills}: Command, Tactics, (optionally) Sway (morale), Investigation (battlefield analysis)
\item \textbf{Talents}: Inspire, Command Presence, Tactical Movement (for allies – see \ref{ch:talents})
\end{itemize}

\textbf{Play Style}:
\begin{itemize}
\item Directing allies in combat (bonus dice, improved Position)
\item Planning and executing tactical maneuvers
\item Leading followers and assets in mass action (see \ref{ch:assets-followers})
\item Using authority to intimidate or persuade
\end{itemize}

\paragraph{Development Path (by Tier)}:
\begin{itemize}
\item \textbf{Tier I}: Raise Presence to 3, Command/Tactics to 2. Take \emph{Command Presence} (4 XP) and \emph{Inspire} (3 XP).
\item \textbf{Tier II}: Take \emph{Tactical Movement} (4 XP) and \emph{Network Builder} (4 XP). Acquire a Cap 2–3 follower (soldier, herald).
\item \textbf{Tier III}: Take \emph{Battlefield Mastery} (6 XP) or \emph{Mystic Bulwark} (8 XP). Build a War Chest (Standard Asset, 8 XP) to fund operations.
\end{itemize}

\paragraph{Tactical Advice}:
\begin{itemize}
\item \textbf{Use \emph{Inspire} to give allies +1 die and +1 Boon.} A well‑timed inspiration can turn a miss into a success.
\item \textbf{Your followers are your strength.} Keep them alive – spend Boons to protect them. A Cap 3 scout is worth more than a single attack.
\item \textbf{Use \emph{Tactical Movement} to reposition without disengaging.} Order allies to flank while you hold the line.
\item \textbf{In mass combat, treat your unit as a high‑Cap follower.} Your Command roll can advance the War Clock faster than any sword.
\end{itemize}

\textbf{Story Hooks}:
\begin{itemize}
\item What army or company did they lead before?
\item What strategic defeat haunts them?
\item Who are their rivals in the art of war?
\item What cause do they fight for — coin, country, or creed?
\end{itemize}

\textbf{Legacy Note}: A loyal lieutenant or sergeant (follower) can inherit your command, leading the company after you. See \ref{ch:legacy-engine}.

\paragraph{Build Blocks.}
\textbf{Starting Build (30 XP).}
\begin{itemize}
\item \textbf{Attributes}: Presence 3 (9), Wits 2 (6), Body 1 (3), Spirit 1 (3) $\rightarrow$ \textbf{21 XP}
\item \textbf{Skills}: Command 2 (4), Tactics 2 (4) $\rightarrow$ \textbf{8 XP}
\item \textbf{Total}: 29 XP (bank 1 XP)
\end{itemize}
\textbf{With Bonds/Complications (34 XP).}
\begin{itemize}
\item Add \textbf{Talent}: Command Presence (4 XP) — +1 die on leadership/intimidation rolls, using banked 1 + 4 = 5 XP; bank 1 XP
\item \textbf{Total}: 33 XP (bank 1 XP)
\end{itemize}

\begin{tcolorbox}[colback=gray!5!white,colframe=gray!70!black,title=Reskinning the Paths: Runekeeper and Invoker as Familiar Concepts,label=box:reskinning]
\index{reskinning}\index{Runekeeper!reskinning}\index{Invoker!reskinning}
    The magic systems in \textit{Fate's Edge} are designed to be reimagined. Neither \textbf{Runekeeper} (Thiasos + Codex) nor \textbf{Invoker} (Patron’s Symbol) is locked into a single narrative identity. Below are common reskins — use the same mechanics, change the fiction.
    
    \subsection*{Runekeeper – The Devoted Agent}
    A Runekeeper has sworn themselves to a single Patron, learning its Rites and bearing its Obligation. This framework can represent:
    
    \begin{itemize}
      \item \textbf{Paladin / Crusader}: Sworn to the \textit{Oath of Flame \& Light} or the \textit{Inquisitor Prime}. Rites become prayers, smites, and consecrations. The Codex is a prayer book or a vow‑etched gauntlet.
      \item \textbf{Warlock / Pact‑Bearer}: Bound to a distant, alien, or dangerous Patron (e.g., Khemesh, Malachai, the Ninth). Rites are granted, not learned. Obligation is the steady pull of a bargain.
      \item \textbf{Cleric / Priest}: Devoted to the Everflame, the Pale Shepherd, or other divine powers. Rites are liturgies. The Patron’s Gift is a blessing or a relic.
      \item \textbf{Dark Knight / Death Knight}: A fallen or morally ambiguous warrior who draws power from a grim Patron (Carrion‑King, Mor’iraath). Melee and Rites mix; Obligation is a slow doom.
    \end{itemize}
    
    \noindent \textit{Mechanical core unchanged:} 1 action Rites, Obligation track, Patron’s Gift imbuement, Push It option.
    
    \subsection*{Invoker – The Professional Occultist}
    An Invoker does not fully commit to a single Patron. Instead, they carry \textbf{Symbols} (often tools, badges, or relics) of several Powers and perform rituals to invoke their Rites. This framework can represent:
    
    \begin{itemize}
      \item \textbf{Exorcist / Witch Hunter}: Carries symbols of several protective Patrons (The Sealed Gate, Oath of Flame, Gallow’s Bell). Rites are wards, banishments, and purification rituals. \emph{Crack the Seal} is an emergency exorcism.
      \item \textbf{Alienist / Scholar of the Unseen}: Studies forbidden Patrons (The Ninth, Khemesh, Mor’iraath) from a safe distance — “safe” being relative. Uses Symbol libraries to invoke their power without full devotion. Rituals are experimental and dangerous.
      \item \textbf{Consulting Occultist / Paranormal Investigator}: Carries a briefcase or satchel of Symbols — a bell from the Mistlands, a sun‑seal from Ecktoria, a hag’s stone from Aelinnel. Hired to solve magical problems. Rituals are filed like work orders; \emph{Crack the Seal} is “burning a favor.”
      \item \textbf{Ritualist / Ceremonial Mage}: Believes that power lies in procedure, not faith. Each Symbol is a “key” — a geometric diagram, a calendar, a consecrated blade. Rites are slow, deliberate, and precise. The Invoker is a technician, not a worshipper.
    \end{itemize}
    
    \noindent \textit{Mechanical core unchanged:} Ritual time, Obligation per invocation, Symbol display reduces cost, \emph{Crack the Seal} for instant effect.
    
    \subsection*{Mixing and Blurring Lines}
    Characters may also blend these frameworks — a Runekeeper who occasionally carries a second Symbol for emergency rites, or an Invoker who one day pledges fully to a Patron and becomes a Runekeeper. The system does not penalize narrative growth; simply retrain Talents during downtime (\S\ref{ch:talents}) when the fiction justifies it.
    
    \noindent \textit{Remember:} The fiction leads. If you want to play a Paladin, describe your Runekeeper’s Rites as prayers. If you want a Vodouisant or Hoodoo practitioner, reflavor Patron’s Symbols as sacred bundles, veves, or spirit vessels. If you want a monster hunter, carry a satchel of monster‑specific Symbols (silver stake, iron bell, cold iron nail). The mechanics stay the same; the story breathes.
    \end{tcolorbox}
    
\section{Creating Your Own Concept}
\label{sec:creating-own-concept}
\index{character concepts!creation}

\subsection*{Start with Narrative}
\begin{itemize}
\item What is your character's background and motivation?
\item What role do they play in their community or society?
\item What relationships are important to them?
\item What goals are they pursuing?
\end{itemize}

\subsection*{Add Mechanical Support}
\begin{itemize}
\item Choose attributes that support your concept (see \ref{sec:core-attributes})
\item Select skills that reflect their training and experience (see \ref{sec:skills})
\item Consider talents that provide unique capabilities (see \ref{ch:talents})
\item Think about assets that represent their resources (see \ref{ch:assets-followers})
\end{itemize}

\subsection*{Consider Group Role}
\begin{itemize}
\item How does your concept complement other party members?
\item What gaps in group capability can you fill?
\item What unique contributions can you make?
\item How will you work with other characters?
\end{itemize}

\subsection*{Plan for Growth}
\begin{itemize}
\item What short-term improvements make sense?
\item What long-term development aligns with your concept?
\item How might your character change over time?
\item What legacy do you want to build? (See \ref{ch:legacy-engine})
\end{itemize}

\begin{tcolorbox}[colback=green!5!white,colframe=green!75!black,title=Character Concept Worksheet,fonttitle=\bfseries]
\index{character worksheet}
\textbf{Narrative Elements:}
\begin{itemize}
\item Concept: \_\_\_\_\_\_\_\_\_\_\_\_\_\_\_\_\_\_\_\_\_\_
\item Traditional Analogues: \_\_\_\_\_\_\_\_\_\_\_\_\_\_\_\_
\item Motivation: \_\_\_\_\_\_\_\_\_\_\_\_\_\_\_\_\_\_\_\_\_\_
\item Background: \_\_\_\_\_\_\_\_\_\_\_\_\_\_\_\_\_\_\_\_\_\_
\item Relationships: \_\_\_\_\_\_\_\_\_\_\_\_\_\_\_\_\_\_\_\_\_\_
\end{itemize}

\textbf{Mechanical Foundation:}
\begin{itemize}
\item Primary Attributes: \_\_\_\_\_\_\_\_\_\_\_\_\_\_\_\_\_\_\_\_\_\_
\item Key Skills: \_\_\_\_\_\_\_\_\_\_\_\_\_\_\_\_\_\_\_\_\_\_
\item Starting Talents: \_\_\_\_\_\_\_\_\_\_\_\_\_\_\_\_\_\_\_\_\_\_
\item Initial Assets: \_\_\_\_\_\_\_\_\_\_\_\_\_\_\_\_\_\_\_\_\_\_
\end{itemize}

\textbf{Development Plan:}
\begin{itemize}
\item Short-term goals: \_\_\_\_\_\_\_\_\_\_\_\_\_\_\_\_\_\_\_\_\_\_
\item Long-term vision: \_\_\_\_\_\_\_\_\_\_\_\_\_\_\_\_\_\_\_\_\_\_
\end{itemize}
\end{tcolorbox}

\section{Final Notes}
\label{sec:concepts-final}
\index{character concepts!final notes}

Remember that these examples are starting points, not limitations. The most interesting characters often combine elements from multiple concepts or create entirely new approaches. Work with your Game Master to ensure your character concept fits the campaign and provides engaging story opportunities.

The best characters are those that you find interesting to play and that contribute to an enjoyable experience for everyone at the table.